\begin{table}[p]
    \centering
    \caption{\\ Bond statistics}
    \scriptsize
    \label{tab:bond_stats}
    
    \begin{threeparttable}
        \begin{tabular}{lccccccc}
            \toprule
                 & N firms & N firms & N bonds & N bonds & Mean $d$ & Median $d$ & $\rho(d_{1930}, d_t)$ \\
            Year & (with bonds) & (with gold bonds) & (all) & (gold) & ($d>0$) & ($d>0$) & \\
            \midrule
            1930 & 157     & 150 & 275 & 250 & 0.57 & 0.55 & 1.00 \\
            1931 & 150     & 148 & 260 & 247 & 0.59 & 0.55 & 0.81 \\
            1932 & 146     & 144 & 251 & 240 & 0.58 & 0.56 & 0.80 \\
            1933 & 134     & 130 & 225 & 212 & 0.59 & 0.59 & 0.71 \\
            1934 & 123     & 114 & 209 & 189 & 0.56 & 0.52 & 0.55 \\
            \bottomrule
        \end{tabular}
    \end{threeparttable}\\
    
    \vspace*{3mm} \justifying \noindent
    \scriptsize{\textit{Notes.} This table examines whether firms strategically adjusted their gold clause exposure prior to the abrogation. The sample is a balanced panel of 157 firms with corporate bonds outstanding at the end of 1930---the year before Britain's departure from the gold standard. For each year, the table reports the number of the initial 157 firms that still have at least one bond outstanding, the number of firms with gold clause bonds, total bond count among the 157 firms, gold clause bond count, mean and median gold clause exposure ($d$, as defined in equation (\ref{eq:d})) among firms with positive exposure, and the correlation between contemporaneous $d$ and its 1930 value.}
\end{table}